\documentclass[10pt,a4paper]{amsart}
\usepackage[margin=19mm]{geometry}
\usepackage{amsmath,amssymb,mathtools}
\usepackage[scr=rsfs,cal=euler]{mathalfa}
\usepackage{xfrac}
\usepackage[unicode,hidelinks,bookmarks=false]{hyperref}
\newcommand{\ie}{i.e.\ }
\newcommand{\cf}{cf.\ }
\newcommand{\resp}{respectively\ }
\newcommand{\Subsection}{\subsection}
\providecommand{\nobreakdash}{\nobreak\hbox{-}}
\setlength{\emergencystretch}{2em}
\multlinegap0pt
\usepackage{bm}
\usepackage{bbm}
\usepackage{dsfont}
\usepackage{xspace}
\usepackage{cancel}
\usepackage{mathtools}
\usepackage{amsthm}
\makeatletter
\@ifundefined{thm}{\newtheorem{thm}{Thm}}{}
\@ifundefined{prop}{\newtheorem{prop}[thm]{Proposition}}{}
\makeatother
\def \al{\alpha}
\def \b {\beta}
\def \ga {\gamma}
\def \h {\mathfrak{h}}
\def\Res{\mathrm{Res}}
\def\ds{\dots}
\title[Finite induction functor for vertex operator algebras]{Finite induction functor for vertex operator algebras}
\author{Jianqi Liu}
\date{Source: arXiv:2503.23632v2 (CC BY 4.0)}
\begin{document}
\maketitle

\noindent\emph{Source and licence.} Finite induction functor for vertex operator algebras. Version arXiv:2503.23632v2 (CC BY 4.0), distributed under the Creative Commons Attribution 4.0 International licence, as recorded in the arXiv licence metadata for this exact version and shown on the abstract page https://arxiv.org/abs/2503.23632v2. Full rights notice: licenses/arxiv-2503.23632-CC-BY-4.0.txt.\par\smallskip

\noindent\textit{Editorial scope.} Counted unit: the Prop whose source label is \texttt{prop4.5} in the article (source0.tex lines 1949--1955, source label \texttt{prop4.5}), delivered together with the article's own complete original proof of it (source0.tex lines 1957--2018, one \texttt{proof} environment of 62 lines). No mathematics is written, repaired or re-derived: statement and proof are the article's own text line for line. Required context retained uncounted: 4.7 (lines 1797--1805); 4.8 (lines 1809--1811). Every definition, display and label of those lines is retained unchanged. Omitted from this package and not redistributed: 1--1796, 1806--1808, 1812--1948, 1956--1956, 2019--3548. Nothing inside a retained excerpt is omitted or altered: every retained body is delivered exactly as the article's lines are written, so no omission receipt of any kind accompanies this package. Citations: the retained excerpts carry no citation command at all, so no bibliography entry of the article is transcribed and no reference is redistributed. Reference adaptation: none was needed and none was made; every reference inside the delivered unit resolves inside it, so no unresolved reference or citation is produced. Printed slips kept literal: no printed word, symbol, hypothesis or number of the counted statement or of its proof was altered. Layout and class: the article's own class is \texttt{amsart} with packages including amsmath, physics, amssymb, graphicx, amsthm, enumitem, tikz, mathrsfs, hyperref, inputenc, tikz-cd, fontenc, txfonts, lipsum; the wrapper is the lane's pinned \texttt{amsart} editorial wrapper at 10pt on A4, which restates the theorem-like environments the retained text needs and adds the article's own macro definitions that the retained text needs (\texttt{\textbackslash Res}, \texttt{\textbackslash al}, \texttt{\textbackslash b}, \texttt{\textbackslash ds}, \texttt{\textbackslash ga}, \texttt{\textbackslash h}), with extra wrapper packages (bm, bbm, dsfont, xspace, cancel, mathtools). The delivered numbering is the unit's own. Assets: no retained span contains a figure environment, a graphics-inclusion command, a PDF-page inclusion, a table or a TikZ picture; no image file is copied into this package, the package declares no asset dependency and no image file is redistributed with it. Licence: version arXiv:2503.23632v2 of this article is distributed under the Creative Commons Attribution 4.0 International licence, as recorded in the arXiv licence metadata for this exact version and shown on the abstract page https://arxiv.org/abs/2503.23632v2. A full rights notice is delivered at licenses/arxiv-2503.23632-CC-BY-4.0.txt. Characters: every retained excerpt is pure ASCII, so the delivered engine, XeLaTeX with its default Latin Modern text font, reports no missing character.
		\begin{equation}\label{4.7}
			\begin{cases}
				&h(-n-2)u+h(-n-1)u,\qquad u\in V_P,\ h\in \mathfrak{h},\ n\geq 0;\\[4pt]
				&\gamma(-1)v+v,\qquad v\in M_{\hat{\mathfrak{h}}}(1,\gamma),\ \gamma\in \{\alpha,-\alpha,\beta,\alpha+\beta\};\\[4pt]
				&\gamma(-1)^2v+\gamma(-1)v,\quad v\in M_{\hat{\mathfrak{h}}}(1,\gamma+\gamma'),\ \gamma,\gamma'\in \{\alpha,-\alpha,\beta,\alpha+\beta\},\gamma+\gamma'\in \{\alpha+\beta,\beta\};\\[4pt]
				&M_{\hat{\mathfrak{h}}}(1,m\alpha+n\beta),\quad m\alpha+n\beta\in (\mathbb{Z}\alpha\oplus\mathbb{Z}_{\geq 0}\beta)\setminus\{0,\alpha,-\alpha,\beta,\alpha+\beta\};\\[4pt]
				&\alpha(-1)^3w-\alpha(-1)w,\qquad w\in M_{\hat{\mathfrak{h}}}(1,0).
			\end{cases}
		\end{equation}

	\begin{equation}\label{4.8}
		\{(\alpha,\beta),(\beta,\alpha),(-\alpha,\alpha+\beta),(\alpha+\beta,-\alpha)\}.
	\end{equation}

	\begin{prop}\label{prop4.5}
		Let $(\ga,\ga')$ be an ordered pair in the set~\eqref{4.8}, and let $n\geq 0$. Then 
		\begin{equation}\label{4.14}
			\Res_z Y(e^{\ga},z)\big(h^1(-n_1)\cdots h^r(-n_r)e^{\ga'}\big)\frac{(1+z)}{z^{2+n}}\in O,
		\end{equation}
		where $r\geq 0$, $h^i\in \h$ for all $i$, and $n_1\geq \cdots \geq n_r\geq 1$.
	\end{prop}

	\begin{proof}
		It is clear from~\eqref{4.7} that $h(-m)O\subset O$ for any $m\geq 1$ and $h\in\h$.  
		We first claim that $L(-1)u+L(0)u\in O$ for all $u\in V_P$.  
		Indeed, if $u\in M_{\hat{\h}}(1,m\al+n\b)$ with $m\al+n\b\in P\setminus\{\al,-\al,\b,\al+\b\}$,  
		then $L(-1)u+L(0)u\in M_{\hat{\h}}(1,m\al+n\b)\subset O$ by~\eqref{4.7}.  
		
		Now let $u=h^1(-n_1)\cdots h^r(-n_r)e^\ga$, where $\ga\in\{\al,-\al,\b,\al+\b\}$,  
		$h^i\in\h$ for all $i$, and $n_1\geq\cdots\geq n_r\geq 1$.  
		Since $L(-1)e^\ga=\ga(-1)e^\ga$, we obtain
		\begin{align*}
			L(-1)u+L(0)u
			&=h^1(-n_1)\cdots h^r(-n_r)(\ga(-1)e^\ga+e^\ga)\\
			&\quad+\sum_{j=1}^r (h^j(-n_j-1)+h^j(-n_j))\,h^1(-n_1)\cdots \widehat{h^j(-n_j)}\cdots h^r(-n_r)e^\ga\\
			&\equiv 0\pmod{O},
		\end{align*}
		by~\eqref{4.7}. Hence $L(-1)u+L(0)u\in O$ for all $u\in V_P$.
		
		Now take $(\ga,\ga')\in\{(\al,\b),(\b,\al),(-\al,\al+\b),(\al+\b,-\al)\}$.  
		Note that $\epsilon(\ga,\ga')=-1$.
		We first use induction on $n\geq 0$ to prove that
		\begin{equation}\label{4.15}
			\Res_z Y(e^{\ga},z)e^{\ga'}\frac{(1+z)}{z^{2+n}}\in O,
		\end{equation}
		which will serve as the base case for induction on $r$ in~\eqref{4.14}.
		
		For $n=0$, we have
		\[
		\Res_z Y(e^\ga,z)e^{\ga'}\frac{(1+z)}{z^{2}}
		= e^\ga_{-2}e^{\ga'}+e^\ga_{-1}e^{\ga'}
		= -\frac{1}{2}\ga(-1)^2e^{\ga+\ga'}-\frac{1}{2}\ga(-1)e^{\ga+\ga'}
		\equiv 0\pmod{O}.
		\]
		Assume~\eqref{4.15} holds for smaller $n$. Then
		\begin{align*}
			&(n+1)(n+2)\big(e^\ga_{-n-3}e^{\ga'}+e^\ga_{-n-2}e^{\ga'}\big)
			=(n+1)(L(-1)e^\ga)_{-n-2}e^{\ga'}+(n+2)(L(-1)e^\ga)_{-n-1}e^{\ga'}\\
			&=(n+1)(\ga(-1)e^\ga)_{-n-2}e^{\ga'}+(n+2)(\ga(-1)e^\ga)_{-n-1}e^{\ga'}\\
			&=(n+1)\!\sum_{j\geq 0}\!\ga(-1-j)e^\ga_{-n-2+j}e^{\ga'}
			+(n+1)\!\sum_{j\geq 0}\! e^\ga_{-n-3-j}\ga(j)e^{\ga'}\\
			&\quad+(n+2)\!\sum_{j\geq 0}\!\ga(-1-j)e^\ga_{-n-1+j}e^{\ga'}
			+(n+2)\!\sum_{j\geq 0}\! e^\ga_{-n-2-j}\ga(j)e^{\ga'}\\
			&=(n+1)\sum_{j\geq 0}\ga(-1-j)e^\ga_{-n-2+j}e^{\ga'}
			+(n+1)(\ga|\ga')e^{\ga}_{-n-3}e^{\ga'}\\
			&\quad+(n+2)\sum_{j\geq 0}\ga(-1-j)e^\ga_{-n-1+j}e^{\ga'}
			+(n+2)(\ga|\ga')e^{\ga}_{-n-2}e^{\ga'}.
		\end{align*}
		Since $(\ga|\ga')=-1$ and $\ga(-2-j)v+\ga(-1-j)v\in O$ for any $j\geq 0$, we have \begin{align*} &(n+1)(n+3)\left(e^\ga_{-n-3}e^{\ga'}+e^\ga_{-n-2}e^{\ga'}\right)\\
			&=(n+1)\ga(-1)e^\ga_{-n-2}e^{\ga'}+(n+1)\sum_{t\geq 0}\ga(-2-t)e^\ga_{-n-1+t}e^{\ga'}\\
			&\quad +(n+2)\sum_{j\geq 0} \ga(-1-j)e^\ga_{-n-1+j}e^{\ga'}-e^\ga_{-n-2}e^{\ga'}\\ &\equiv (n+1)\ga(-1)e^\ga_{-n-2}e^{\ga'}+\sum_{j\geq 0} \ga(-1-j) e^\ga_{-n-1+j}e^{\ga'}-e^\ga_{-n-2}e^{\ga'}\pmod{O}\\ &\equiv (n+1)\ga(-1)e^\ga_{-n-2}e^{\ga'}+\sum_{j\geq 0} (-1)^j\ga(-1) e^\ga_{-n-1+j}e^{\ga'}-e^\ga_{-n-2}e^{\ga'}\pmod{O}. 
		\end{align*}
		
		Since $\ga(-1)O\subset O$, then by the induction hypothesis we have \begin{align*} \ga(-1)e^\ga_{-n-1+j}e^{\ga'}\equiv (-1)e^\ga_{-n-1+j-1}e^{\ga'}\equiv \ds\equiv (-1)^j \ga(-1)e^\ga_{-n-1}e^{\ga'}\pmod{O}, \end{align*} for any $0\leq j\leq n$, and $e^\ga_{-n-2}e^{\ga'}\equiv \ds \equiv (-1)^{n+1}e^\ga_{-1}e^{\ga'}=(-1)^{n+1}\epsilon(\ga,\ga')\ga(-1)e^{\ga+\ga'}\pmod{O}$. Moreover, we have $e^{\ga}_{0}e^{\ga'}=\epsilon(\ga,\ga')e^{\ga+\ga'}$, and $e^\ga_{m}e^{\ga'}=0$ for $m\geq 1$.
		It follows that \begin{align*} &(n+1)\ga(-1)e^\ga_{-n-2}e^{\ga'}+\sum_{j=0}^{n+1} (-1)^j\ga(-1) e^\ga_{-n-1+j}e^{\ga'}-e^\ga_{-n-2}e^{\ga'}\\ &\equiv (n+1) \ga(-1)e^{\ga}_{-n-2}e^{\ga'}+(n+1)\ga(-1)e^{\ga}_{-n-1}e^{\ga'}+ (-1)^{n+1} \ga(-1)e^{\ga}_{0}e^{\ga'}-e^\ga_{-n-2}e^{\ga'}\\ &\equiv 0+ (-1)^{n+1} \ga(-1)\epsilon(\ga,\ga')e^{\ga+\ga'}-(-1)^{n+1}\epsilon(\ga,\ga')\ga(-1)e^{\ga+\ga'}\\ &\equiv 0\pmod{O}. \end{align*}
		Hence $(n+1)(n+3)\big(e^\ga_{-n-3}e^{\ga'}+e^\ga_{-n-2}e^{\ga'}\big)\in O$, 
		completing the proof of~\eqref{4.15}.
		
		Finally, we use induction on the length $r$ to prove~\eqref{4.14}.  
		The base case $r=0$ follows from~\eqref{4.15}.  
		Assume~\eqref{4.14} holds for smaller $r\geq 1$. Then \begin{align*} &\Res_z Y(e^\ga,z)h^1(-n_1)\ds h^r(-n_r)e^{\ga'} \frac{(1+z)}{z^{2+n}}\\ &= h^1(-n_1)\ds h^r(-n_r) \Res_z Y(e^\ga,z)e^{\ga'} \frac{(1+z)}{z^{2+n}}\\ &-\sum_{j=1}^r (h^j|\ga) \Res_z h^1(-n_1)\ds h^{j-1}(-n_{j-1}) Y(e^\ga,z)h^{j+1}(-n_{j+1})\ds h^r(-n_r)e^{\ga'} \frac{(1+z)}{z^{2+n+n_j}}\\ &\equiv 0\pmod{O}, \end{align*}
		where the last congruence follows from the induction hypothesis and the fact that 
		$h(-m)O\subset O$ for all $h\in\h$ and $m\geq 1$.
	\end{proof}

\end{document}
